\section{有效形式形变与 Grothendieck 存在性定理}\label{src:C.5}
\begin{flushright}\itshape
如果你理解了某件事, 就会理解它是显然的.\\
\normalfont Israel Gelfand
\end{flushright}
Grothendieck 存在性定理是证明形式形变有效的一种强有力的工具. 它有时也称为形式 GAGA, 因为它类似于 Serre 的 GAGA 定理 \cite{Ser56}: 对 $\mathbb C$ 上的固有概形 $X$ 及其解析化 $X^{\mathrm{an}}$, 将凝聚层 $F$ 映到其解析化 $F^{\mathrm{an}}$ 的自然函子 $\operatorname{Coh}(X)\to\operatorname{Coh}(X^{\mathrm{an}})$ 是等价. 两者的证明遵循非常相似的策略.

\subsection{有效形式形变}\label{src:C.5.1-subsection}
\begin{definition}\label{src:C.5.1}
设 $\mathcal X$ 为 $\Sch/\Bbbk$ 上的预叠 (或函子). 取 $x_0\in\mathcal X(\Bbbk)$, 并考虑 $x_0$ 在一个剩余域为 $\Bbbk$ 的完备 Noether 局部 $\Bbbk$-代数上的形式形变 $\{x_n\}$ (更准确地说, 是注记 \bookref{C.4.4} 中定义的 $x_0$ 处形变叠 $\mathcal X_{x_0}$ 的形式形变). 若存在对象 $\widehat x\in\mathcal X(R)$ 及相容的同构 $x_n\xrightarrow\sim\widehat x|_{\Spec R/\mathfrak m^{n+1}}$, 则称 $\{x_n\}$ \newterm{有效的}{effective}.
\end{definition}
形式形变 $(R,\{x_n\})$ 有效, 当且仅当它属于自然函子 $\mathcal X(R)\to\varprojlim\mathcal X(R/\mathfrak m^n)$ 的本质像. 换言之, 存在使下图交换的虚线箭头:
\[
\begin{tikzcd}[column sep=small]
\Spec R/\mathfrak m\arrow[r,hook]\arrow[drrrr,bend right=12,"x_0"'] & \Spec R/\mathfrak m^2\arrow[r,hook]\arrow[drrr,bend right=8,"x_1"'] & \Spec R/\mathfrak m^3\arrow[r,hook]\arrow[drr,"x_2"'] & \cdots\arrow[r,hook] & \Spec R\arrow[d,dashed,"\widehat x"]\\
&&&&\mathcal X
\end{tikzcd}
\]
\begin{example}\label{src:C.5.2}
若反变函子 $F:\Sch/\Bbbk\to\mathrm{Sets}$ 由 $\Bbbk$ 上的概形 $X$ 表示, 则每个形式形变 $(R,\{x_n\})$ 都有效. 事实上, $x_n$ 对应于像为 $x\in X(\Bbbk)$ 的态射 $\Spec R/\mathfrak m^{n+1}\to X$, 因而对应于 $\Bbbk$-代数同态 $\phi_n:\widehat{\mathcal O}_{X,x}\to R/\mathfrak m^{n+1}$. 对 $\phi_n$ 取逆像, 得到局部同态 $\widehat{\mathcal O}_{X,x}\to R$, 它又定义了延拓 $\{x_n\}$ 的态射 $\widehat x:\Spec R\to X$.

更一般地, 若 $\mathcal X$ 是 $\Bbbk$ 上的代数叠, 则每个形式形变都有效. 事实上, 存在光滑呈示 $U\to\mathcal X$ 及 $x_0\in\mathcal X(\Bbbk)$ 的提升 $u\in U(\Bbbk)$. 应用光滑性的无穷小提升判别法 (\bookref{4.7.1}), 可归纳构造提升
\[
\begin{tikzcd}
\Spec R/\mathfrak m^n\arrow[r]\arrow[d,hook]&U\arrow[d]\\
\Spec R/\mathfrak m^{n+1}\arrow[r,"x_n"']\arrow[ur,dashed]&\mathcal X.
\end{tikzcd}
\]
因为 $U$ 是概形, 映射 $\Spec R/\mathfrak m^n\to U$ 延拓为映射 $\Spec R\to U$, 复合 $\Spec R\to U\to\mathcal X$ 使形式形变 $\widehat x=\{x_n\}$ 有效化. 因为 $\mathcal X$ 的对角可表, $x_n$ 的相容自同构 $\alpha_n$ 延拓为 $\widehat X$ 的唯一自同构, 即函子
\[
\mathcal X(R)\to\varprojlim\mathcal X(R/\mathfrak m^{n+1})
\]
是范畴等价.
\end{example}

\subsection{Grothendieck 存在性定理}\label{src:C.5.2-subsection}
下述定理常用于 $(R,\mathfrak m)$ 为完备局部环的情形.
\begin{theorem}[Grothendieck 存在性定理]\label{src:C.5.3}
设 $X$ 为 Noether 环 $A$ 上的固有概形, 且 $A$ 关于理想 $\mathfrak m\subseteq A$ 完备. 令 $X_n:=X\times_RR/\mathfrak m^{n+1}$. 函子
\begin{equation}\label{src:C.5.4}
\operatorname{Coh}(X)\to\varprojlim\operatorname{Coh}(X_n),\qquad E\mapsto\{E/\mathfrak m^{n+1}E\}
\end{equation}
是范畴等价.
\end{theorem}
\bookeqref{C.5.4} 的本质满性意味着下图可延拓:
\[
\begin{tikzcd}[column sep=small]
E_0\arrow[d,no head]&E_1\arrow[d,no head]&E_2\arrow[d,no head]&&E\arrow[d,no head]\\
X_0\arrow[r,hook]\arrow[d]&X_1\arrow[r,hook]\arrow[d]&X_2\arrow[r,hook]\arrow[d]&\cdots\arrow[r,hook]&X\arrow[d]\\
\Spec R/\mathfrak m\arrow[r,hook]&\Spec R/\mathfrak m^2\arrow[r,hook]&\Spec R/\mathfrak m^3\arrow[r,hook]&\cdots\arrow[r,hook]&\Spec R.
\end{tikzcd}
\]
用形式概形的语言, 令 $\widehat X=X\times_{\Spec R}\operatorname{Spf} R$ 为 $X$ 的 $\mathfrak m$-进完备化, 则 Grothendieck 存在性定理断言, 由 $E\mapsto\widehat E$ 定义的函子 $\operatorname{Coh}(X)\to\operatorname{Coh}(\widehat X)$ 为等价.
\begin{remark}[范畴的极限]\label{src:C.5.5}
$\varprojlim\operatorname{Coh}(X_n)$ 的对象是由凝聚 $\mathcal O_{X_n}$-模 $E_n$ 组成的序列 $\{E_n\}$, 配备同构 $E_{n+1}/\mathfrak m^{n+1}E_{n+1}\xrightarrow\sim E_n$. 态射 $\phi=\{\phi_n\}:\{E_n\}\to\{F_n\}$ 是相容的 $\mathcal O_{X_n}$-模同态 $\phi_n:E_n\to F_n$. 若 $X\cong\Spec R$, 则 $\varprojlim\operatorname{Coh}(X_n)$ 等价于有限 $A$-模的范畴.

范畴 $\varprojlim\operatorname{Coh}(X_n)$ 是阿贝尔范畴, 且 \bookeqref{C.5.4} 是正合函子. 虽然其中映射 $\{\phi_n\}:\{E_n\}\to\{F_n\}$ 的余核就是 $\{\operatorname{coker}\phi_n\}$, 核却更为微妙. Artin--Rees 引理 (\bookref{B.5.4}) 蕴含, 对每个 $m$, $\ker(\alpha_l)\to E_m$ 的像在 $l\gg m$ 时稳定为子层 $E_m'\subseteq E_m$; 对每个 $n$, 商 $E_m'/\mathfrak m^{n+1}E_m'$ 在 $m\gg n$ 时稳定为 $X_n$ 上的凝聚层 $K_n$. $\{\phi_n\}$ 的核为 $\{K_n\}$ (通常不同于 $\{\ker\phi_n\}$). 见 \cite[Tag 0EHN]{SP}.
\end{remark}
\begin{proof}
\bookeqref{C.5.4} 的全忠实性等价于双射
\[
\Ext^0_{\mathcal O_X}(E,F)\xrightarrow\sim\varprojlim_n\Ext^0_{\mathcal O_{X_n}}(E|_{X_n},F|_{X_n}),
\]
这是形式函数定理的推论 (见习题 \bookref{A.5.5}). 剩下只需证明本质满性.

\textit{射影情形:} 关键断言是, 对每个 $\{E_n\}\in\operatorname{Coh}(X_n)$, 存在整数 $m,r$ 及相容的满射 $\mathcal O_{X_n}(-m)^{\oplus r}\to E_n$. 考虑有限型 $\mathcal O_{X_0}$-代数 $\mathcal A:=\bigoplus_{i\geq0}\mathfrak m^i\mathcal O_{X_0}/\mathfrak m^{i+1}\mathcal O_{X_0}$ 及有限生成 $\mathcal O_{\mathcal A}$-模 $\mathcal G:=\bigoplus_{i\geq0}\mathfrak m^iE_i$ (注意, 当 $\{E_n\}=\{E/\mathfrak m^{n+1}\}$, 其中 $E$ 为 $X$ 上的凝聚层时, $\mathcal G$ 就是伴随分次模 $\bigoplus_{i\geq0}\mathfrak m^iE/\mathfrak m^{i+1}E$). 将 $\mathcal G$ 看作 $\Spec_{X_0}\mathcal A$ 上的凝聚层, 对射影态射 $\Spec_{X_0}\mathcal A\to\Spec\bigoplus_{i\geq0}\mathfrak m^i/\mathfrak m^{i+1}$ 应用 Serre 消灭定理 \cite[定理 II.5.2]{Har77}, 得到整数 $m_0$, 使得对所有 $m\geq m_0$, $\mathrm H^1(X_0,\mathcal G(m))=0$. 因而对所有 $n$, $\mathrm H^1(X_0,\mathfrak m^{n+1}E_n(m))=0$, 从而
\begin{equation}\label{src:C.5.6}
\mathrm H^0(X_0,E_{n+1}(m))\to\mathrm H^0(X_0,E_n(m))
\end{equation}
满射. 必要时增大 $m_0$, 可保证 $E_0(m)$ 由截面 $s_{0,1},\ldots,s_{0,r}$ 整体生成. 由 \bookeqref{C.5.6} 的满射性, 可找到截面 $s_{0,i}$ 的相容提升 $s_{n,i}\in\mathrm H^0(X_0,E_n(m))$. 这给出相容的映射 $\mathcal O_{X_n}^{\oplus r}\to E_n(m)$, 由 Nakayama 引理, 每个映射均满射.

断言得证后, 令 $\{K_n\}\in\varprojlim\operatorname{Coh}(X_n)$ 为 $\{\phi_n:\mathcal O_{X_n}(-m)^{\oplus r}\to E_n\}$ 的核, 如注记 \bookref{C.5.5} 所述. 再对 $\{K_n\}$ 应用上述断言, 得到相容的右正合列
\[
\mathcal O_{X_n}(-m')^{\oplus r'}\xrightarrow{\alpha_n}\mathcal O_{X_n}(-m)^{\oplus r}\to E_n\to0.
\]
由全忠实性, 态射 $\{\alpha_n\}$ 由态射 $\widetilde\alpha:\mathcal O_X(-m')^{\oplus r'}\to\mathcal O_X(-m)^{\oplus r}$ 诱导, 因而 $E:=\operatorname{coker}(\widetilde\alpha)$ 是 $X$ 上的凝聚层, 满足 $\varprojlim\operatorname{Coh}(X_n)$ 中的 $\{E/\mathfrak m^{n+1}E\}\cong\{E_n\}$.

\textit{固有情形:} 由 Chow 引理 \cite[习题 II.4.10]{Har77}, 存在射影态射 $g:X'\to X$, 它在稠密开子集 $U\subseteq X$ 上为同构, 且 $X'$ 在 $\Spec R$ 上射影. 给定 $\{E_n\}\in\operatorname{Coh}(X_n)$, 考虑拉回 $\{g^*E_n\}\in\operatorname{Coh}(X_n')$, 其中 $X_n':=X'\times_RR/\mathfrak m^{n+1}$. 因为 $X'$ 在 $R$ 上射影, 存在 $X'$ 上的凝聚层 $E'$ 及 $\varprojlim\operatorname{Coh}(X_n')$ 中的同构
\[
\{\beta_n\}:\{g^*E_n\}\xrightarrow\sim\{E'/\mathfrak m^{n+1}E'\}.
\]
由上同调的有限性 (\bookref{A.5.3}), $g_*E'$ 凝聚. 从概念上说, 接下来的论证十分直接: 伴随映射 $E_n\to g_*g^*E_n$ 诱导 $\varprojlim\operatorname{Coh}(X_n)$ 中的正合列
\begin{equation}\label{src:C.5.7}
\begin{aligned}
0\to\{K_n\}\to\{E_n\}\xrightarrow{\{\alpha_n\}}&\{(g_*E')/\mathfrak m^{n+1}(g_*E')\}\\
&\xrightarrow{\{q_n\}}\{Q_n\}\to0,
\end{aligned}
\end{equation}
其中 $\ker(\alpha)=\{K_n\}$ 及 $\operatorname{coker}(\alpha)=\{Q_n\}$ 支撑于 $X\setminus U$, 即每个 $K_n,Q_n$ 都支撑于 $X\setminus U$. 由 Noether 归纳, 可假设存在 $X$ 上的凝聚层 $K,Q$ 及同构 $\{K_n\}\cong\{K/\mathfrak m^{n+1}K\}$ 和 $\{Q_n\}\cong\{Q/\mathfrak m^{n+1}Q\}$. 由 \bookeqref{C.5.4} 的全忠实性与正合性, 映射 $\{q_n\}$ 由满射 $g_*E'\to Q$ 诱导, 定义 $F:=\ker(g_*E'\to Q)$. 因为
\[
\Ext^i_{\mathcal O_X}(K,F)\xrightarrow\sim\varprojlim_n\Ext^i_{\mathcal O_{X_n}}(K_n,F/\mathfrak m^{n+1}F),
\]
存在扩张 $0\to K\to E\to F\to0$, 给出 $X$ 上的凝聚层 $E$, 使得 $\{E_n\}\cong(E/\mathfrak m^{n+1}E)$.

然而, \bookeqref{C.5.7} 中 $\{\alpha_n\}$ 的存在性需要一些工作. 此处形式概形的形式体系很有用: 考虑 $\operatorname{Spf} R$ 上的形式概形映射 $\widehat g:\widehat X'\to\widehat X$, $\widehat X$ 上的凝聚层 $\widehat E=\varprojlim E_n$, 以及 $\operatorname{Coh}(\widehat X)$ 中的伴随态射 $\widehat E\to\widehat g_*\widehat g^*\widehat E$, 应用形式函数定理的一个版本 \cite[III$_1$.4.1.5]{EGA}, 有时称为 ``比较定理'', 得到等同
\[
\mathrm R^i\widehat g_*(\widehat g^*\widehat E)\cong\varprojlim\mathrm R^ng_*(g^*E_n)\cong(\mathrm R^ng_*E')^{\widehat{\ }}.
\]
我们给出更直接的论证. 首先证明存在唯一的映射 $\alpha_n$, 补全下图:
\[
\begin{tikzcd}
E_n\arrow[r,"c_n"]\arrow[d,dashed,"\alpha_n"']&g_*g^*E_n\arrow[d,"g_*\beta_n","\sim"']\\
(g_*E')/\mathfrak m^{n+1}(g_*E')\arrow[r,"d_n"']&g_*(E'/\mathfrak m^{n+1}E'),
\end{tikzcd}
\]
其中 $c_n,d_n$ 为自然映射. 由唯一性, $\{\alpha_n\}$ 的存在性是 $X$ 上的局部问题, 故可假设 $X=\Spec B$. 因为 $B\to\widehat B:=\varprojlim B/\mathfrak m^{n+1}B$ 平坦, 且图中所有凝聚层均被 $\mathfrak m$ 的某个幂零化, 平坦基变换 (\bookref{A.2.12}) 进一步将问题归结为 $B$ 关于 $\mathfrak mB$ 完备的情形. 此时 $E:=\varprojlim E_n$ 对应于 $X$ 上的凝聚层, 在 $\varprojlim\operatorname{Coh}(X_n)$ 中映到 $\{E_n\}$. 对 $X'\to\Spec B$ 及 $E'$ 应用形式函数定理 (\bookref{A.5.4}), 得到 $g_*E'\cong\varprojlim g_*(E'/\mathfrak m^{n+1}E')$, 表明 $\varprojlim d_n$ 为同构. 因此复合 $\alpha:=(\varprojlim d_n)^{-1}\circ\varprojlim(g_*\beta_n\circ c_n)$ 定义映射 $E\to g_*E'$, 诱导所需的映射 $\alpha_n$.

现在证明 \bookeqref{C.5.7} 的核与余核支撑于 $X\setminus U$. 因为 $g^*E\cong\varprojlim g^*E_n$, 同构 $\{\beta_n\}:\{g^*E_n\}\to\{E'/\mathfrak m^{n+1}E'\}$ 来自同构 $g^*E\xrightarrow\sim E'$, 从而 $\{\alpha_n\}:\{E_n\}\to\{(g_*E')/\mathfrak m^{n+1}(g_*E')\}$ 来自 $E\to g_*g^*E\xrightarrow\sim g_*E'$. 因为 $g$ 在 $U$ 上为同构, 伴随映射 $E\to g_*g^*E$ 在 $U$ 上为同构, 因而 $\ker\alpha_n$ 和 $\operatorname{coker}\alpha_n$ 支撑于 $X\setminus U$. 另见 \cite[III$_1$.5.1.4]{EGA}, \cite[定理 $8.4.2$]{Ill05} 和 \cite[Tag 088E]{SP}.
\end{proof}

\begin{corollary}\label{src:C.5.8}
设 $(R,\mathfrak m)$ 为完备 Noether 局部环, $X_n\to\Spec R/\mathfrak m^{n+1}$ 为一列固有态射, 满足 $X_n\times_{R/\mathfrak m^{n+1}}R/\mathfrak m^n\cong X_{n-1}$. 若 $L_n$ 是 $X_n$ 上的相容线丛序列, 且 $L_0$ 丰沛, 则存在射影态射 $X\to\Spec R$, $X$ 上的丰沛线丛 $L$, 及相容的同构 $X_n\cong X\times_RR/\mathfrak m^{n+1}$ 和 $L_n\xrightarrow\sim L|_{X_n}$.
\end{corollary}
换言之, 下述笛卡尔图可延拓, 其中 $X$ 在 $R$ 上射影:
\[
\begin{tikzcd}[column sep=small]
X_0\arrow[r,hook]\arrow[d]&X_1\arrow[r,hook]\arrow[d]&X_2\arrow[r,hook]\arrow[d]&\cdots\arrow[r,hook,dashed]&X\arrow[d,dashed]\\
\Spec R/\mathfrak m\arrow[r,hook]&\Spec R/\mathfrak m^2\arrow[r,hook]&\Spec R/\mathfrak m^3\arrow[r,hook]&\cdots\arrow[r,hook]&\Spec R.
\end{tikzcd}
\]
称 $X_0$ 的形式形变 $\{X_n\to\Spec R/\mathfrak m^{n+1}\}$ 有效.\footnote{有时也称为可代数化, 但我们将这个术语留给有限型 $\Bbbk$-概形上的形变, 例如 Artin 代数化 (定理 \bookref{C.6.8}) 的结论.}
\begin{proof}
令 $\Bbbk=R/\mathfrak m$. 考虑有限生成分次 $\Bbbk$-代数 $B=\bigoplus\mathfrak m^n/\mathfrak m^{n+1}$ 及拟凝聚分次 $\mathcal O_{X_0}$-代数 $\mathcal A=B\otimes_{\Bbbk}\mathcal O_{X_0}$. 对 $\Spec_{X_0}\mathcal A$ 及丰沛线丛 $L_0\otimes_{\mathcal O_{X_0}}\mathcal O_{X_0'}$ 应用 Serre 消灭定理, 得到 $d_0$, 使得 $d\geq d_0$ 时 $\mathrm H^1(X_0,\mathcal A\otimes L_0^{\otimes d})=0$. 必要时增大 $d_0$, 可假设 $L_0^{d_0}$ 极丰沛. 取 $\mathrm H^0(X_0,L_0^{\otimes d_0})$ 的截面 $s_{0,0},\ldots,s_{0,N}$, 定义闭浸入 $X_0\hookrightarrow\mathbb P^N$. 有正合列
\[
0\to\mathfrak m^n\mathcal O_{X_{n+1}}/\mathfrak m^{n+1}\mathcal O_{X_{n+1}}\to\mathcal O_{X_{n+1}}\to\mathcal O_{X_n}\to0.
\]
与 $L_{n+1}^{\otimes d}$ 张量得到短正合列
\[
0\to(\mathfrak m^n\mathcal O_{X_{n+1}}/\mathfrak m^{n+1}\mathcal O_{X_{n+1}})\otimes L_0^{\otimes d}\to L_{n+1}^{\otimes d}\to L_n^{\otimes d}\to0,
\]
这里使用了 $(\mathfrak m^n\mathcal O_{X_{n+1}}/\mathfrak m^{n+1}\mathcal O_{X_{n+1}})$ 支撑于 $X_0$, 以及 $m\leq n$ 时的等同 $L_{n+1}\otimes\mathcal O_{X_m}\cong L_m$. $\mathrm H^1(X_0,\mathcal A\otimes L_0^{\otimes d})$ 的消失蕴含, 可将 $s_{0,0},\ldots,s_{0,N}$ 归纳提升为 $\mathrm H^0(X_n,L_n^{\otimes d})$ 的相容截面 $s_{n,0},\ldots,s_{n,N}$. 由 Nakayama 引理, 诱导态射 $X_n\hookrightarrow\mathbb P^N_{R/\mathfrak m^{n+1}}$ 为闭浸入, 给出交换图
\[
\begin{tikzcd}[column sep=small,row sep=large]
&\mathbb P^N\arrow[rr]\arrow[dd]&&\mathbb P^N_{R/\mathfrak m_R^2}\arrow[rr]\arrow[dd]&&\cdots\arrow[rr]&&\mathbb P_R^N\arrow[dd]\\
X_0\arrow[ur,hook,"\text{闭}"]\arrow[rr]\arrow[dr]&&X_1\arrow[ur,hook,"\text{闭}"]\arrow[rr]\arrow[dr]&&\cdots\arrow[rr,dashed]&&X\arrow[ur,hook,dashed,"\text{闭}"]\arrow[dr,dashed]\\
&\Spec\Bbbk\arrow[rr,hook]&&\Spec R/\mathfrak m^2\arrow[rr,hook]&&\cdots\arrow[rr,hook]&&\Spec R.
\end{tikzcd}
\]
Grothendieck 存在性定理 (\bookref{C.5.3}) 给出等价 $\operatorname{Coh}(\mathbb P_R^N)\to\varprojlim\operatorname{Coh}(\mathbb P^N_{R/\mathfrak m^{n+1}})$. 本质满性给出 $\mathbb P_R^N$ 上延拓 $\{\mathcal O_{X_n}\}$ 的凝聚层 $E$, 全忠实性给出延拓 $\mathcal O_{\mathbb P^N_{R/\mathfrak m^{n+1}}}\to\mathcal O_{X_n}$ 的满射 $\mathcal O_{\mathbb P_R^N}\to E$. 取 $X\subseteq\mathbb P_R^N$ 为 $\ker(\mathcal O_{\mathbb P_R^N}\to E)$ 定义的闭子概形. 另见 \cite[III.5.4.5]{EGA}, \cite[定理 $8.4.10$]{Ill05} 及 \cite[Tag 089A]{SP}.
\end{proof}
\begin{remark}\label{src:C.5.9}
另假设每个 $X_n$ 在 $R/\mathfrak m^{n+1}$ 上平坦. 若只给定 $X_0$ 上的丰沛线丛 $L_0$ (而未给定线丛 $L_n$), 则将 $L_{n-1}$ 形变为 $L_n$ 的障碍是元素 $\mathrm{ob}_{L_{n-1}}\in\mathrm H^2(X,\mathcal O_X\otimes_{\Bbbk}\mathfrak m^n)$ (命题 \bookref{C.2.11}). 若这些上同调群消失 (例如 $X$ 的维数为 $1$), 则存在相容的延拓 $L_n$, 从而形式形变 $\{X_n\to\Spec R/\mathfrak m^{n+1}\}$ 有效.

若 $L_0$ 的形变 $L_n$ 不存在, 则形式形变未必有效. 例如, 存在射影 K3 曲面 $(X_0,L_0)$ 及非射影的一阶形变 $X_1\to\Spec\Bbbk[\epsilon]$ (所以 $L_0$ 无法形变到 $X_1$), 并存在无效的形式形变; 见 \cite[例 $21.2.1$]{Har10}, \cite[断言 $3.5$]{Sta06} 和 \cite[Tag 0D1Q]{SP}. 类似地, 阿贝尔簇的形式形变也可能无效. 注意, 对抽象 K3 曲面或阿贝尔簇的模问题, Rim--Schlessinger 判别法 (\bookref{C.4.6}) 可用于构造万有形式形变, 但缺少有效性意味着相应的叠不是代数叠 (例 \bookref{C.5.2}).
\end{remark}
我们将 Grothendieck 存在性定理应用于 Hilbert 函子 $\Hilb(X)$, 所有曲线的叠 $\mathcal M_g^{\mathrm{all}}$ 及凝聚层叠 $\Coh(X)$, 它们均如命题 \bookref{C.4.14} 中定义在 $\Sch/\Bbbk$ 上.
\begin{proposition}\label{src:C.5.10}
对 $\Hilb(X)$, $\mathcal M_g^{\mathrm{all}}$ 和 $\Coh(X)$, 每个形式形变都有效. 特别地, 存在有效的极小万有形式形变.
\end{proposition}
\begin{proof}
对 $\Hilb(X)$, 设 $\{Z_n\subseteq X_{R/\mathfrak m^{n+1}}\}$ 为形式形变. Grothendieck 存在性定理 (\bookref{C.5.3}) 蕴含, 存在 $X_R$ 上延拓结构层 $\{\mathcal O_{Z_n}\}$ 的凝聚层 $E$, 并且存在延拓 $\{\mathcal O_{X_n}\to\mathcal O_{Z_n}\}$ 的满射 $\mathcal O_{X_R}\to E$. 由 $\ker(\mathcal O_{X_R}\to E)$ 定义的子概形 $Z\subseteq X_R$ 使形式形变有效化.

对 $\mathcal M_g^{\mathrm{all}}$, 设 $\{C_n\to\Spec R/\mathfrak m^{n+1}\}$ 为形式形变. 因为 $C_0$ 是域上的固有曲线, 它是射影的. 取 $C_0$ 上的丰沛线丛 $L_0$. 将 $C_n$ 上的线丛 $L_n$ 形变到 $C_{n+1}$ 上的线丛 $L_{n+1}$ 的障碍是元素 $\mathrm{ob}_{L_{n-1}}\in\mathrm H^2(X,\mathcal O_X\otimes_{\Bbbk}\mathfrak m^n)$ (命题 \bookref{C.2.11}). 因为 $\dim C=1$, 此上同调群为零, 因而存在线丛的相容族 $\{L_n\}$. 于是可应用推论 \bookref{C.5.8}.

对 $\Coh(X)$, 形式形变的有效性直接来自 Grothendieck 存在性定理 (\bookref{C.5.3}). 最后一个断言来自极小万有形式形变的存在性 (命题 \bookref{C.4.14}).
\end{proof}
\begin{exercise}\label{src:C.5.11}
设 $X,Y$ 为完备 Noether 局部环 $R$ 的谱 $S=\Spec R$ 上的固有概形. 记 $X_n,Y_n$ 为 $X,Y$ 在 $S_n=\Spec R/\mathfrak m_R^{n+1}$ 上的限制. 证明 $S_n$ 上的相容态射序列 $f_n:X_n\to Y_n$ 延拓为唯一态射 $f:X\to Y$.
\end{exercise}
\subsection{提升到特征 0}\label{src:C.5.3-subsection}
形变理论的一个引人注目的应用是将 $\operatorname{char}(\Bbbk)=p$ 的域 $\Bbbk$ 上的概形 $X_0$ ``提升'' 到特征 $0$. 若存在特征 $0$ 的完备 Noether 局部环 $(R,\mathfrak m)$, 满足 $R/\mathfrak m=\Bbbk$, 及光滑概形 $X\to\Spec R$, 满足 $X_0\cong X\times_R\Bbbk$, 则称 $X_0$ \newterm{可提升}{liftable}到特征 $0$.\footnote{此定义有若干变体, 例如已经给定剩余域为 $\Bbbk$ 的完备离散赋值环 $R$ 的情形.} 可以希望在 $X$ 上使用特征 $0$ 的工具 (例如 Hodge 理论), 推出 $X_0$ 的性质. 提升 $X_0$ 的策略是归纳地将 $X_0$ 形变为 $R/\mathfrak m^{n+1}$ 上的概形 $X_n$, 再应用 Grothendieck 存在性定理使形式形变有效化.

光滑曲线可提升, 因为形变曲线和丰沛线丛的障碍均消失. Serre 给出了不可提升的射影三维簇的例子 (见 \cite[定理 $22.4$]{Har10}), Mumford 将其推广为不可提升的射影曲面 (见 \cite[推论 $8.6.7$]{Ill05}). 另一方面, Mumford 证明主极化阿贝尔簇可提升 \cite{Mum69}, Deligne 证明 K3 曲面可提升 \cite{Del81}. 这些例子相当有趣, 因为在这两种情形中, 形式形变均未必有效 (见注记 \bookref{C.5.9}).
